% Generated draft: requires mathematical and source review.
\documentclass[11pt]{article}
\usepackage[margin=25mm]{geometry}
\usepackage{fontspec}
\setmainfont{Latin Modern Roman}
\usepackage{amsmath,amssymb,mathtools,unicode-math}
\setmathfont{Latin Modern Math}
\usepackage{enumitem,needspace,xurl,hyperref}
\setlength{\parindent}{0pt}
\setlength{\parskip}{5pt}
\setlength{\emergencystretch}{3em}
\newcommand{\N}{\mathbb{N}}
\newcommand{\Z}{\mathbb{Z}}
\newcommand{\Q}{\mathbb{Q}}
\newcommand{\R}{\mathbb{R}}
\newcommand{\C}{\mathbb{C}}
\newcommand{\F}{\mathbb{F}}
\newcommand{\A}{\mathbb{A}}
\newcommand{\PP}{\mathbb{P}}
\newcommand{\E}{\mathbb{E}}
\newcommand{\eps}{\varepsilon}
\newcommand{\dd}{\,\mathrm{d}}
\begin{document}
\section*{UnsolvedMath ID: 30005412}
\textbf{Source problem code:} catalogue entry
\subsection{Definitions and mathematical statement}
Let \((R,\mathfrak m)\) be a commutative Noetherian local Gorenstein ring, using the standard convention for this terminology, and let \(I\subsetneq R\) be a perfect ideal of codimension \(3\). Thus
\[
\operatorname{pd}_R(R/I)=\operatorname{grade}_R(I)=3.
\]
Write a minimal free resolution as
\[
0\longrightarrow F_3\xrightarrow{d_3}F_2\xrightarrow{d_2}F_1\xrightarrow{d_1}R\longrightarrow R/I\longrightarrow0.
\]
Its ranks give the resolution format, recorded here as \((1,\operatorname{rank}F_1,\operatorname{rank}F_2,\operatorname{rank}F_3)\). This notation does not specify what the source means by a \emph{smaller} format.

For ideals \(J,K\subsetneq R\) of codimension \(3\), a direct complete-intersection link is given by an ideal \(C=(a_1,a_2,a_3)\subseteq J\cap K\), generated by an \(R\)-regular sequence, such that
\[
J=C:K,\qquad K=C:J,
\qquad C:J=\{r\in R:rJ\subseteq C\}.
\]
The ideal \(I\) is \emph{LICCI} if there is a finite chain of such links
\[
I=I_0\sim I_1\sim\cdots\sim I_s
\]
ending at a complete-intersection ideal; a chain of length zero is allowed.

\textbf{Conjecture 2, as supplied.} For every such \(R\) and \(I\), the proposed criterion is
\[
I\text{ is LICCI}
\quad\Longleftrightarrow\quad
\text{there exists a higher structure map in }W(d_1)
\text{ having a unit entry}.
\]
A unit entry means a coefficient in \(R^\times=R\setminus\mathfrak m\). Here \(W(d_i)\) retains the source's notation for representations containing higher structure maps; it is not being identified with the ordinary differential \(d_i\). The supplied record does not define these representations, the allowed higher structure maps, their construction from a resolution, or the choices involved. Consequently, the right-hand condition remains source-dependent and cannot be made self-contained from this record alone.

The accompanying conjectural assertion asks whether the existence of a unit in \(W(d_2)\) and \(W(d_3)\) is equivalent to the existence of an ideal \(I'\) linked to \(I\) whose quotient has a resolution of smaller format. The record does not specify whether the two unit conditions are separate criteria or a joint condition, whether linkage here means one direct link or a finite chain, or the comparison of formats. No particular interpretation of these unresolved points is imposed in this draft.

\textbf{Supplied contextual qualifications.} The excerpt also asserts that one can construct a uniform family of generic examples of perfect ideals whose quotients have resolutions of ADE formats, and says these examples are \emph{expected} to provide generic models for the resolutions. It describes three viewpoints: bigradings of Lie algebras of type \(T_{p,q,r}\), with uniform descriptions of the differentials, referring to its reference [10]; defining ideals of affine pieces of Schubert varieties in certain homogeneous spaces of groups of type \(T_{p,q,r}\), referring to [12]; and partial Jacobian ideals of sporadic invariants of prehomogeneous representations in the Sato--Kimura list, also referring to [12]. The listed invariants are \(\Delta_4\), of degree \(4\) for \(\operatorname{SL}(6)\), \(\Delta_7\), of degree \(7\) for \(\operatorname{SL}(7)\), and \(\Delta_{16}\), of degree \(16\) for \(\operatorname{SL}(8)\). Their ambient representation formulas are corrupted in the supplied extraction and are not reconstructed here. The excerpt says that type \(D_n\) recovers the Buchsbaum--Eisenbud generic Gorenstein ideals of codimension \(3\), citing [4], and Anne Brown's almost complete intersections of codimension \(3\), citing [2]. Definitions of ADE formats and the bibliographic identities of [2], [4], [10], and [12] are not supplied.

Finally, the excerpt introduces a Theorem 3 concerning perfect codimension-three ideals over local Gorenstein rings whose quotient resolutions have ADE format, mentioning joint work with L. Guerrieri and X. Ni, but ends immediately after the word ``Then.'' Its conclusion is missing, so this fragment supplies no complete theorem statement.
\subsection{Short English statement}
The conjecture asks whether every perfect codimension-three ideal over a local Gorenstein ring belongs to the complete-intersection linkage class exactly when a higher structure map in \(W(d_1)\) has a unit coefficient. Related unit conditions involving \(W(d_2)\) and \(W(d_3)\) are proposed to detect linkage to smaller resolution formats. The higher-map constructions and the meaning of the latter criterion are not fully specified in the supplied record. The supplied dated curation labels the problem partially solved; this draft does not independently verify that assessment.
\subsection{Sources}
\begin{itemize}
\item Supplied local UnsolvedMath registry record, permanent ID \(30005412\), title ``Higher Structure Map Criterion for LICCI Ideals,'' exact source problem code \texttt{catalogue entry}, in the Oberwolfach Reports Open Problems collection; proposed year \(2023\). No individual proposer is identified in the record.
\item \emph{Resolutions in Local Algebra and Singularity Theory} (2023), Oberwolfach Reports, pp.\ 327--395, DOI \texttt{10.4171/owr/2023/6}: \url{https://doi.org/10.4171/owr/2023/6}. The supplied excerpt identifies the target as Conjecture 2 and ends with a truncated Theorem 3 near the printed page marker 357.
\item Supplied literature reference: Guerrieri--Ni--Weyman, ``higher structure maps'': \url{https://arxiv.org/abs/2208.05934}. The record provides this attribution and link, without a full bibliographic citation.
\item Supplied literature reference: ``OWR-2023-6, Conjecture 2'': \url{https://oa.tib.eu/renate/server/api/core/bitstreams/11773ea5-b1e5-4e9e-b650-808e9d9dc944/content}.
\item Literature assessment dated exactly \texttt{2026-08-22}: supplied dated curation, not independently verified. It labels the problem ``partially solved'' and reports that the LICCI criterion is proved for the smallest previously unknown ADE format and related linkage reductions, while the uniform if-and-only-if conjecture for all codimension-three perfect ideals remains open. The record itself requests independent verification of time-sensitive status and cited claims.
\item The registry labels the statement \texttt{corrected\_verified} and reports that the official Oberwolfach report was checked directly. These are supplied provenance claims, not verification performed for this draft. Internal references [2], [4], [10], and [12] occur in the excerpt, but their bibliographic entries are absent from the supplied record.
\end{itemize}
\end{document}
